% Generated draft: requires mathematical and source review.
\documentclass[11pt]{article}
\usepackage[margin=25mm]{geometry}
\usepackage{fontspec}
\setmainfont{Latin Modern Roman}
\usepackage{amsmath,amssymb,mathtools,unicode-math}
\setmathfont{Latin Modern Math}
\usepackage{enumitem,needspace,xurl,hyperref}
\setlength{\parindent}{0pt}
\setlength{\parskip}{5pt}
\setlength{\emergencystretch}{3em}
\newcommand{\N}{\mathbb{N}}
\newcommand{\Z}{\mathbb{Z}}
\newcommand{\Q}{\mathbb{Q}}
\newcommand{\R}{\mathbb{R}}
\newcommand{\C}{\mathbb{C}}
\newcommand{\F}{\mathbb{F}}
\newcommand{\A}{\mathbb{A}}
\newcommand{\PP}{\mathbb{P}}
\newcommand{\E}{\mathbb{E}}
\newcommand{\eps}{\varepsilon}
\newcommand{\dd}{\,\mathrm{d}}
\begin{document}
\section*{UnsolvedMath ID: 30000696}
\textbf{Source problem code:} catalogue entry
\subsection{Definitions and mathematical statement}
Let $\sigma_1$ be the variance-one semicircle law with density $(2\pi)^{-1}\sqrt{4-x^2}$ on $[-2,2]$, and let $\sigma_s$ be the law of $\sqrt{s}W$ for $W\sim\sigma_1$, with $\sigma_0=\delta_0$. A classical variance-scale mixture is
\[
 \mu_\nu(B)=\int_{[0,\infty)}\sigma_s(B)\,\nu(ds).
\]
In a standard independent-clock interpretation of subordinated matrix Brownian motion, $\nu=\nu_t$ is the law of a nonnegative subordinator at time $t$, and the mixture describes the corresponding averaged scalar law in the specified spectral limit. The source does not provide the exact clock, normalization, dimension/limit or averaging construction; these must be matched before deciding the class of admissible mixtures. It does not ask about every arbitrary mixing measure without the subordination qualification.

For probability laws on $\R$, free additive convolution $\boxplus$ is the distribution of the sum of freely independent self-adjoint variables. Free infinite divisibility means
\[
 \forall m\geq1\ \exists\text{ probability law }\mu_m:\qquad
                         \mu=\mu_m^{\boxplus m}.
\]
The first question asks whether the scale mixtures arising from the report's construction satisfy this property. The second explicitly asks it for the classical Gaussian law, for example $N(0,1)$.

This is \emph{free additive} infinite divisibility, not ordinary convolution infinite divisibility or free multiplicative infinite divisibility. Classical scale mixing also must not be identified without proof with a free multiplicative convolution involving the Wigner law: independent scalar multiplication and multiplication of free noncommutative variables generally define different operations. An arbitrary-clock empirical spectral limit may be random before averaging, which is another reason to retain the exact construction.

The literature triage checked 21 August 2026 reports the classical Gaussian freely infinitely divisible, and substantial criteria for free multiplicative Wigner mixtures, while retaining an incomplete general classification for the report's mixtures. The positive Gaussian result answers that named subquestion in the supplied assessment; it does not prove the first assertion for all admissible clocks, and the multiplicative-mixture criteria need their exact comparison hypotheses. These are dated supplied findings, not independently verified theorem/status claims.
\subsection{Short English statement}
Are the specified subordinated semicircle scale mixtures freely additively infinitely divisible? The source also singles out the Gaussian law, which the supplied assessment reports as settled positively. Classical scale mixing and free multiplicative convolution must be distinguished before transferring special-case criteria.
\subsection{Sources}
\begin{itemize}
\item Local UnsolvedMath registry, permanent ID 30000696, exact source code catalogue entry; \emph{Free Infinite Divisibility of Semicircle Scale Mixtures}. Proposed year 2007, status partially\_solved; corrected\_verified statement preserves both the subordination-mixture and Gaussian questions.
\item \emph{Mini-Workshop: L\'evy Processes and Related Topics in Modelling (2007)}, Oberwolfach Reports, DOI 10.4171/owr/2007/07, pp. 417--446, \url{https://doi.org/10.4171/owr/2007/07}.
\item Supplied literature triage checked 2026-08-21: ``original problem report'', \url{https://doi.org/10.4171/owr/2007/07}; ``normal distribution is freely infinitely divisible'', \url{https://arxiv.org/abs/0910.4263}; ``free multiplicative Wigner mixtures'', \url{https://arxiv.org/abs/0910.1199}. This dated OWR-derived curation requests independent verification of time-sensitive status and cited claims.
\end{itemize}
\end{document}
